\cotefile{91}
\scanpage{19}
marqués de $(\unsure{+|\star|\ast})$, dont les images dans le coin droit du bas sont les mêmes, on cherche un élément du coin gauche du haut qui donne ces éléments marqués. On note qu'il suffit pour ceci que $P(u_0,S)\ne\varnothing$, car alors le diagramme $(1)$ est isomorphe à $(2)$, or $(1)$ est \uline{cartésien}, grâce au fait que $M_{v_0}\to M_{u_0}$ est surjectif. Or l'obstruction à trouver un élément dans $P(u_0,S)$ est \deleted{\unsure{dans}}

\[
\partial(u_0)=\xi\in\operatorname{Ext}^{1}_{\unsure{\mathcal O_{P_0}|\underline{\mathcal O}_{P_0}}}
(P_0;I_{\Gamma_0},\mathcal O_{\Gamma_0}\otimes_{S_0}J)
\]

% slanted

\margin{\deleted{\unsure{sous réserve|sauf erreur} qu'il ne faut pas remplacer $P_0$ par $P_1$, \unsure{évident que|évident par}}}

où $P_0=X_0\times_{S_0}Y_0$, $\deleted{\unsure{\xi|\Gamma}}\Gamma_0$ est le sous-préschéma de $P_0$ graphe de $u_0$, et où $I_{\Gamma_0}$ est l'Idéal de $P_0$ qui le définit. C'est un module sur $A_0$, qui s'annule quand on \deleted{\unsure{tensorise par $A_{0\eta}$}}\inserted{localise en $\eta$}, donc Si $\pi$ est un paramètre de $A$, il existe un $n\geq 0$ tel que $\pi^n\xi=0$. Utilisant la suite exacte

\[
0\to J\xrightarrow{\pi^{\unsure{n|n+1}}}J\to J_n\to 0 \tag*{$(\times)$}
\]

et la suite correspondante

\[
0\to\mathcal O_{\Gamma_0}\otimes_{S_0}J\xrightarrow{\pi^n}\underline{\mathcal O}_{\Gamma_0}\otimes_{S_0}J\to\mathcal O_{\unsure{\Gamma_0|\Gamma_s}}\otimes_{S_0}J_n\to0
\]

et la suite exacte des $\operatorname{Ext}^{i}_{\mathcal O_{P_0}}(P_0;I_{\Gamma_0},-)$, on trouve que $\xi$ provient d'un élément

\[
\begin{aligned}
&\xi'\in\operatorname{Hom}_{\mathcal O_{P_0}}(P_0;I_{\Gamma_0},\underline{\mathcal O}_{\Gamma_0}\otimes_{S_0}J_n)\\
&\hspace{4em}\bmod\operatorname{Im}\operatorname{Hom}_{\mathcal O_{P_0}}(P_0;I_{\Gamma_0},\mathcal O_{\Gamma_0}\otimes_{S_0}J]
\end{aligned}
\]

% archive sheet number: 18
